\section{Open Questions}

Five substantive questions this replication did not close. Each is grounded in a
specific gap identified during the paper re-read and has concrete next steps.

\begin{enumerate}

\item \textbf{Dose-dependence of miscount bias.}
\emph{Does the microscope-conditioned miscount bias vary systematically with
delivered dose, or is it stable across the 0.25--4\,Gy range the paper spans?}
Basis: the re-pass results (Claims 8--10) are reported at fixed 1 or 2\,Gy
conditions. Fig 1 uses multiple doses but the miscount summaries aggregate
across dose within each of the 24 microscope configs, so a dose-conditional
bias curve is not exhibited in either the paper's headline claims or in
\texttt{ALL\_CLAIMS\_SUMMARY.json}. If miscount grows with dose (because
clustering grows with local DSB density), a single-dose bias correction is
unsafe for low-dose radiobiology.
\emph{Next steps:} Fit a \%-miscount(dose) regression per microscope config
using the existing 24 parquets; report slope, intercept, and dose-conditional
95\% CI. If slope is nonzero, publish a dose-conditional correction table.
Requires no new simulations --- data is already on disk.

\item \textbf{Generalization across cell lines / chromatin state.}
\emph{How robust is the miscount characterization to cell-line differences
that affect the effective DSB clustering the LoG counter sees?} Basis: Geant4-DNA
generates spatial DSB distributions from radiation transport in a generic
nuclear geometry; the paper does not vary chromatin compaction, heterochromatin
fraction, or nuclear volume across cell lines. Because Fig 8 shows clustering is
the dominant driver of DSB-under-counting, a chromatin-dense cell line
(lymphocytes vs fibroblasts vs iPSCs) would plausibly shift the entire miscount
curve, but this is not tested. \emph{Next steps:} Rerun Geant4-DNA with a
compact-nucleus geometry (or scale coordinate density by a chromatin-compaction
factor $c \in \{0.5, 1.0, 2.0\}$) and repeat the Airyscan-x63 pipeline; compare
miscount($c$) curves. Requires the blocked Python-3.11 environment (F6).

\item \textbf{Bias propagation into repair-kinetics fits.}
\emph{Is the bi-exponential repair-kinetics fit ($a_1{=}0.711, a_2{=}0.289,
\tau_1{=}1.54\,\text{h}, \tau_2{=}10\,\text{h}$) still identifiable when the input
DSB counts are replaced by microscope-observed foci counts?} Basis: Claim 13
confirms the bi-exponential curve is uniformly enforced across radiation types
in the simulated actual-DSB stream. But biology users fit the same
bi-exponential to \emph{observed} foci counts, not to ground truth. If the
microscope-dependent bias is dose- or time-dependent, $\tau_1$ and $\tau_2$
estimated from foci counts will be biased in a way this paper does not report.
\emph{Next steps:} Fit the bi-exponential to counted-foci curves per microscope
config from the same 24 parquets; report
$(\tau_{1,\text{obs}}/\tau_{1,\text{true}},\ \tau_{2,\text{obs}}/\tau_{2,\text{true}})$
as a function of magnification. If ratios deviate from 1, publish a per-microscope
correction for repair-kinetics inference. Trivial with existing data.

\item \textbf{Extrapolation to modalities outside Table 1.}
\emph{Do the paper's magnification and voxel-size trends generalize to
modalities the paper does not cover (STED, expansion microscopy, light-sheet),
or is the 24-config sweep sufficient to define a global correction law?} Basis:
Table 1 lists 24 configurations (wide-field, confocal, Airyscan) with voxel
sizes down to $\sim0.028{\times}0.028{\times}0.10\,\mu\text{m}$. Modalities that
reach sub-diffraction voxels (STED at $\sim50\,\text{nm}$ lateral, ExM at
$\sim20\,\text{nm}$ effective) are not simulated. The Fig 5
Spearman $r=-0.78$ is derived from a bounded voxel-size range and may saturate or
invert outside it. \emph{Next steps:} Add STED and ExM PSF configurations to the
pipeline; rerun for one dose $\times$ radiation. If the Fig 5 trend extrapolates,
publish an extended correction law; if it saturates, publish a modality-specific
correction. Requires blocked pipeline (F6).

\item \textbf{Optimality of the CD\_200nm clustering metric.}
\emph{Does CD\_200nm fully capture the counting-relevant clustering, or would a
resolution-matched CD$_r$ (with $r$ matching the microscope PSF FWHM) explain
more variance in miscount?} Basis: Claim 12 confirms CD\_200nm predicts DSB
under-count ($r=-0.088$) but the effect size is modest despite very high
significance from large $n$. 200\,nm is chosen a-priori and is comparable to the
Airyscan-x63 PSF FWHM; for lower-resolution configs (x10, x20) the
counting-relevant clustering radius should be larger. A microscope-conditional
clustering metric might explain more variance and give a cleaner mechanistic
story. \emph{Next steps:} Compute CD$_r$ for $r \in \{100, 200, 400, 800, 1600\}$\,nm
on the Airyscan-x63 rows already on disk; report per-$r$ Spearman vs miscount
and identify $r^*$ that maximizes $|\text{Spearman}|$. If $r^*>200\,$nm, Fig 8
should be re-plotted with $r^*$. Uses existing data only.

\end{enumerate}
